% ===========================================================================
%  evnx 0.5.0 — User Guide
%  Build:  pdflatex -interaction=nonstopmode main.tex   (run twice for TOC)
% ===========================================================================
\documentclass[aspectratio=169,10pt]{beamer}

\newcommand{\capturedir}{../captures}
% ---------------------------------------------------------------------------
%  evnx presentation preamble — shared by all three decks
%
%  Engine: pdfLaTeX (also builds under XeLaTeX / LuaLaTeX).
%  Packages: beamer, listings, xcolor, amssymb, newunicodechar, booktabs,
%            fancyvrb, hyperref — all in texlive-latex-recommended + beamer.
%
%  \capturedir must be defined by the including document BEFORE \input of
%  this file, e.g.  \newcommand{\capturedir}{../captures}
% ---------------------------------------------------------------------------

\usetheme{Madrid}
\usecolortheme{seahorse}
\usefonttheme{professionalfonts}
\setbeamertemplate{navigation symbols}{}
\setbeamertemplate{caption}[numbered]
\setbeamerfont{frametitle}{size=\large}
\setbeamerfont{title}{size=\LARGE,series=\bfseries}

\usepackage[T1]{fontenc}
\usepackage[utf8]{inputenc}
\usepackage{lmodern}
\usepackage{amssymb}
\usepackage{booktabs}
\usepackage{xcolor}
\usepackage{listings}
\usepackage{fancyvrb}
\usepackage{newunicodechar}
\usepackage{tikz}
\usetikzlibrary{arrows.meta,positioning,shapes.geometric,fit,backgrounds}

% ---- palette ---------------------------------------------------------------
\definecolor{evnxInk}{HTML}{1F2933}
\definecolor{evnxAccent}{HTML}{2F6F4E}
\definecolor{evnxWarn}{HTML}{B7791F}
\definecolor{evnxErr}{HTML}{B03A2E}
\definecolor{evnxOk}{HTML}{2F6F4E}
\definecolor{evnxMuted}{HTML}{6B7280}
\definecolor{evnxTermBg}{HTML}{F5F5F2}
\definecolor{evnxRule}{HTML}{D1D5DB}

\setbeamercolor{structure}{fg=evnxAccent}
\setbeamercolor{alerted text}{fg=evnxErr}

% ---- glyph mapping for body text ------------------------------------------
% evnx emits exactly these 11 non-ASCII characters (enumerated from the
% captured logs, not guessed).
\newcommand{\evnxCheck}{\textcolor{evnxOk}{\ensuremath{\checkmark}}}
\newcommand{\evnxCross}{\textcolor{evnxErr}{\ensuremath{\times}}}
\newcommand{\evnxWarnSign}{\textcolor{evnxWarn}{\textbf{!}}}

\newunicodechar{✓}{\evnxCheck}
\newunicodechar{✗}{\evnxCross}
\newunicodechar{⚠}{\evnxWarnSign}
\newunicodechar{→}{\ensuremath{\rightarrow}}
\newunicodechar{↔}{\ensuremath{\leftrightarrow}}
\newunicodechar{─}{\textendash}
\newunicodechar{—}{---}
\newunicodechar{·}{\textperiodcentered}
\newunicodechar{…}{\ldots}
\newunicodechar{•}{\textbullet}
\newunicodechar{️}{}                % U+FE0F variation selector — invisible

% ---- terminal listing style ------------------------------------------------
% `literate` is required: newunicodechar does not apply inside listings.
\lstdefinestyle{evnxterm}{
  basicstyle=\ttfamily\scriptsize,
  backgroundcolor=\color{evnxTermBg},
  frame=single,
  rulecolor=\color{evnxRule},
  framesep=4pt,
  breaklines=true,
  breakatwhitespace=false,
  columns=fullflexible,
  keepspaces=true,
  showstringspaces=false,
  upquote=true,
  extendedchars=true,
  literate=
    {✓}{{\textcolor{evnxOk}{$\checkmark$}}}1
    {✗}{{\textcolor{evnxErr}{$\times$}}}1
    {⚠}{{\textcolor{evnxWarn}{\textbf{!}}}}1
    {→}{{$\rightarrow$}}1
    {↔}{{$\leftrightarrow$}}1
    {─}{{-}}1
    {—}{{---}}1
    {·}{{$\cdot$}}1
    {…}{{\ldots}}1
    {•}{{\textbullet}}1
    {️}{{}}0
    {é}{{\'e}}1
}

\lstdefinestyle{evnxcode}{
  style=evnxterm,
  basicstyle=\ttfamily\scriptsize,
  backgroundcolor=\color{white},
}

\lstset{style=evnxterm}

% ---- helper macros ---------------------------------------------------------
% \capture{file}        — include a real captured run verbatim
% \capturepart{f}{a}{b} — include lines a..b of a captured run
\newcommand{\capture}[1]{\lstinputlisting[style=evnxterm]{\capturedir/#1}}
\newcommand{\capturepart}[3]{%
  \lstinputlisting[style=evnxterm,firstline=#2,lastline=#3]{\capturedir/#1}}

% a labelled "real output" box
\newcommand{\realrun}[1]{%
  \begin{center}\scriptsize\textcolor{evnxMuted}{real run \textendash\ #1}\end{center}}

\newcommand{\cmd}[1]{\texttt{\textbf{#1}}}
\newcommand{\flag}[1]{\texttt{#1}}
\newcommand{\file}[1]{\texttt{#1}}

% exit-code chips
\newcommand{\exitzero}{\textcolor{evnxOk}{\texttt{0}}}
\newcommand{\exitone}{\textcolor{evnxWarn}{\texttt{1}}}
\newcommand{\exittwo}{\textcolor{evnxErr}{\texttt{2}}}

% a standard "verified against" footline for factual slides
\newcommand{\verifiedon}{\tiny\textcolor{evnxMuted}{verified against evnx 0.5.0 @ 9cfe75b}}

\usepackage[colorlinks=true,linkcolor=evnxAccent,urlcolor=evnxAccent]{hyperref}


\title{evnx \texttt{0.5.0}}
\subtitle{A complete command guide}
\author{}
\date{}

\begin{document}

\begin{frame}[plain]
  \titlepage
  \begin{center}\footnotesize
  \textcolor{evnxMuted}{Every output in this deck is a real run, captured from
  \texttt{evnx 0.5.0}, tag \texttt{v0.5.0} (\texttt{85a52c3}).}
  \end{center}
\end{frame}

% ---------------------------------------------------------------------------
\begin{frame}{What evnx is}
  A single static binary that manages \file{.env} files: validation, secret
  scanning, format conversion, and zero-knowledge encrypted sync.

  \vspace{0.8em}
  \begin{block}{The problem it solves}
    Every other config artefact has tooling — schemas, linters, type checking,
    CI. \file{.env} has a text file, a \file{.gitignore} line, and hope. It also
    holds the highest-value secrets in the system.
  \end{block}

  \vspace{0.6em}
  \begin{tabular}{@{}ll@{}}
    Binary size & 8.5\,MB, no runtime required \\
    Language    & Rust (edition 2021, MSRV 1.85) \\
    Licence     & MIT \\
    Install     & 9 channels: curl, brew, scoop, winget, npm, pip, cargo, Action, Docker \\
    Released    & tag \texttt{v0.5.0} = \texttt{85a52c3} \\
  \end{tabular}
\end{frame}

\begin{frame}[fragile]{The command surface}
  \capture{01-help-root.txt}
\end{frame}

\begin{frame}[fragile]{One caveat on installing}
  Every prebuilt channel ships \textbf{all 17 commands}. \texttt{cargo install}
  is the exception --- it compiles from source with \texttt{default = []}, which
  is 11 commands.

  \vspace{0.5em}
\begin{lstlisting}[style=evnxcode,language=bash]
cargo install evnx                    # 11 commands
cargo install evnx --all-features     # all 17
cargo install evnx --features cloud   # + auth, vault, cloud
\end{lstlisting}

  \vspace{0.3em}
  {\footnotesize New in v0.5.0: a build that is missing features says so in
  \texttt{--help}, and names the flag that would add them. The prebuilt binaries
  never show that block.}
\end{frame}

\begin{frame}{The exit-code contract}
  \begin{center}
  \begin{tabular}{@{}cll@{}}
    \toprule
    \textbf{Code} & \textbf{Meaning} & \textbf{Example} \\
    \midrule
    \exitzero & Clean                & no secrets, valid, in sync \\
    \exitone  & A verdict, and it is bad & secrets found, drifted, invalid \\
    \exittwo  & \textbf{No verdict}  & file missing, regex broken, no TTY \\
    \bottomrule
  \end{tabular}
  \end{center}

  \vspace{0.8em}
  \begin{block}{The sentence that makes it legible}
    \texttt{No verdict: evnx did not finish. This is not a clean result.}
  \end{block}

  \vspace{0.5em}
  \begin{alertblock}{Gate on non-zero, never on \texttt{== 1}}
    A pipeline testing \texttt{\$? -eq 1} treats a crash as a pass. That is how
    a broken scanner silently reports a clean repository for six months.
  \end{alertblock}

  {\footnotesize \cmd{backup} and \cmd{restore} extend this with codes 3--6;
  see their section.}
\end{frame}

\begin{frame}[fragile]{Output streams — why redirection is safe}
  \begin{center}
  \begin{tabular}{@{}ll@{}}
    \toprule
    \textbf{stdout} & \textbf{stderr} \\
    \midrule
    JSON, SARIF, GitHub annotations & human headers and trailers \\
    converted formats & \texttt{Docs:} links \\
    generated completions & \texttt{Validation failed: 3 errors} \\
    \bottomrule
  \end{tabular}
  \end{center}

  \vspace{0.8em}
\begin{lstlisting}[style=evnxcode,language=bash]
evnx validate --format json 2>/dev/null | jq '.issues'
evnx scan --format sarif  > results.sarif     # clean SARIF
evnx convert --to json    > config.json       # clean JSON
\end{lstlisting}

  {\footnotesize Verified for every machine format — the human trailer never
  contaminates a redirect.}
\end{frame}

\begin{frame}[fragile]{The five-command core}
\begin{lstlisting}[style=evnxcode,language=bash]
evnx init --from-source   # scaffold from what the code actually reads
evnx scan                 # secrets that look real        (exit 1 = found)
evnx validate             # missing, placeholder, weak, boolean traps
evnx sync --check         # is .env.example still honest?  (exit 1 = drifted)
evnx template --strict    # render a config, refuse to write a broken one
\end{lstlisting}

  \vspace{0.6em}
  All five are non-interactive, all five have a machine format, and all five
  respect the 0/1/2 contract.
\end{frame}

\begin{frame}{Contents}
  \tableofcontents[hideallsubsections]
\end{frame}

% ---------------------------------------------------------------------------
%  Per-command sections. Each file is self-contained and also builds alone
%  via user-guide/standalone/<cmd>.tex
% ---------------------------------------------------------------------------
% ===========================================================================
%  evnx init
% ===========================================================================
\section{evnx init}

\begin{frame}{\cmd{evnx init} — start a project}
  \textbf{Purpose.} Create \file{.env}, \file{.env.example} and a correct
  \file{.gitignore} for a project that has none.

  \vspace{0.6em}
  \begin{block}{Synopsis}
    \texttt{evnx init [OPTIONS]}
  \end{block}

  \vspace{0.4em}
  \begin{tabular}{@{}ll@{}}
    \toprule
    \textbf{Flag} & \textbf{What it does} \\
    \midrule
    \flag{--detect}          & Read the project's \emph{manifests} and scaffold from them \\
    \flag{--from-source}     & Read the project's \emph{code} and scaffold from that \\
    \flag{--with <LIST>}     & Name components directly: \texttt{nextjs,postgresql} \\
    \flag{--blueprint <ID>}  & A named stack, e.g. \texttt{t3\_modern} \\
    \flag{--list-components} & Print everything \flag{--with} accepts, then exit \\
    \flag{-y, --yes}         & Non-interactive; blank unless a stack is named \\
    \flag{--force}           & Replace files that already exist \\
    \flag{--path <PATH>}     & Where to write (default \texttt{.}) \\
    \bottomrule
  \end{tabular}

  \vspace{0.5em}
  {\footnotesize \textbf{Effectively run-once.} A second run refuses rather than
  overwriting a \file{.env.example} you may have hand-maintained — exit \exittwo,
  with \flag{--force} named as the way through.}
\end{frame}

\begin{frame}[fragile]{\cmd{init --detect} — scaffold from the manifest}
  \capture{10-init-detect.txt}
  \realrun{a directory containing only \file{package.json} with \texttt{next}}
\end{frame}

\begin{frame}[fragile]{What \flag{--detect} wrote}
  \capture{10-init-detect-example.txt}

  {\footnotesize Note the \file{.gitignore} line it generated:
  \texttt{.env*, !.env.example, !.env.sample, !.env.template}. The negation is
  the part people get wrong by hand.}
\end{frame}

\begin{frame}[fragile]{\cmd{init --from-source} — scaffold from the code}
  Reads \texttt{process.env.X}, \texttt{os.getenv("X")},
  \texttt{os.environ["X"]} and \texttt{std::env::var("X")}.

  \capture{11-init-from-source.txt}

  {\footnotesize It \emph{refuses} rather than guessing: no terminal and no
  \flag{--yes} is exit \exittwo, and it names the exact command to proceed.}
\end{frame}

\begin{frame}[fragile]{\cmd{init --from-source --yes} — accepted}
  \capture{11b-init-from-source-yes.txt}

  \begin{block}{Why this is the best mode}
    It is derived from what the code \emph{actually reads}, so it answers
    ``is my \file{.env.example} complete?'' from evidence rather than from a
    dependency list. Output is deterministic and each entry carries a
    \texttt{file:line} citation.
  \end{block}
\end{frame}

\begin{frame}[fragile]{\cmd{init --with} — name the stack}
  \capture{12-init-with.txt}
  \realrun{\texttt{evnx init --with nextjs,postgresql,stripe -y}}

  {\footnotesize Discover the vocabulary with \cmd{evnx init --list-components}.}
\end{frame}

\begin{frame}[fragile]{The generated \file{.env}}
  \capturepart{12-init-with-env.txt}{1}{16}

  {\footnotesize \textbf{Changed in v0.5.0:} \file{.env} is written with
  \emph{active assignments}. It previously wrote every line commented as
  \texttt{\# TODO:}, which meant \cmd{evnx validate} — init's own suggested next
  step — reported every variable missing.}
\end{frame}

\begin{frame}[fragile]{When to use which mode}
  \begin{tabular}{@{}lll@{}}
    \toprule
    \textbf{Situation} & \textbf{Use} & \textbf{Why} \\
    \midrule
    Existing codebase      & \flag{--from-source} & Evidence from the code itself \\
    Fresh scaffold         & \flag{--detect}      & Manifest already declares the stack \\
    Known stack, scripted  & \flag{--with}        & Explicit, reproducible, CI-safe \\
    Greenfield, exploring  & (interactive)        & Guided walk \\
    \bottomrule
  \end{tabular}

  \vspace{1em}
  \begin{block}{Re-running is refused, not silently destructive}
    \capture{15-init-rerun.txt}
  \end{block}
\end{frame}

\section{evnx add}

\begin{frame}{\cmd{evnx add} — extend an existing \file{.env}}
  \begin{block}{Synopsis}\texttt{evnx add <SUBCOMMAND> [OPTIONS]}\end{block}

  \begin{tabular}{@{}ll@{}}
    \toprule
    \texttt{add service <NAME>}   & \texttt{postgresql}, \texttt{redis}, \texttt{stripe}\ldots \\
    \texttt{add framework <NAME>} & needs \flag{-l, --language} \\
    \texttt{add blueprint <ID>}   & a whole stack, without overwriting \\
    \texttt{add custom}           & interactive \\
    \midrule
    \flag{-p, --path <PATH>}      & Target directory \\
    \flag{-y, --yes}              & Skip confirmation \\
    \bottomrule
  \end{tabular}

  \vspace{0.5em}
  {\footnotesize Discover the vocabulary with \cmd{evnx init --list-components}.}
\end{frame}

\begin{frame}[fragile]{A real run}
  \capture{45-add-service.txt}

  {\footnotesize It never overwrites a value that is already there, and it
  cross-checks \file{.gitignore} — \cmd{add} is the command most likely to
  introduce a secret-shaped variable into a trackable file.}
\end{frame}

\begin{frame}{Use cases}
  \begin{itemize}
    \item \textbf{Adding Redis to a running project}\\
      {\footnotesize\texttt{evnx add service redis -y}}
    \item \textbf{A Django app inside a monorepo}\\
      {\footnotesize\texttt{evnx add framework --language python django -p services/web}}
    \item \textbf{A whole stack at once, safely}\\
      {\footnotesize\texttt{evnx add blueprint t3\_modern}}
  \end{itemize}

  \vspace{0.8em}
  \begin{block}{\cmd{init} vs \cmd{add}}
    \cmd{init} is for a project with no \file{.env} at all and refuses to
    overwrite one. \cmd{add} is the supported way to extend afterwards.
  \end{block}

  \vspace{0.4em}
  {\footnotesize An unknown component is exit \exittwo{} and names
  \cmd{evnx init --list-components} as the way to find the right one.}
\end{frame}

\section{evnx validate}

\begin{frame}{\cmd{evnx validate} — is this \file{.env} fit to run?}
  \textbf{Purpose.} Compare \file{.env} against \file{.env.example} and, when
  present, the \texttt{[vars.*]} contract in \file{.evnx.toml}.

  \begin{block}{Synopsis}\texttt{evnx validate [OPTIONS]}\end{block}

  \begin{tabular}{@{}ll@{}}
    \toprule
    \textbf{Flag} & \textbf{What it does} \\
    \midrule
    \flag{--env <PATH>}        & Validate this file instead of \file{.env} \\
    \flag{--env-name <NAME>}   & \texttt{production} $\rightarrow$ \file{.env.production} \\
    \flag{--example <PATH>}    & Compare against this template \\
    \flag{--strict}            & Also report variables absent from the template \\
    \flag{--fix}               & Repair what can be repaired safely \\
    \flag{--format <FORMAT>}   & \texttt{pretty} (default), \texttt{json}, \texttt{github} \\
    \flag{--exit-zero}         & Always exit \exitzero{} (advisory mode) \\
    \flag{--ignore <TYPES>}    & Comma-separated issue types to suppress \\
    \flag{--validate-formats}  & Also check value shapes: url, port, email \\
    \bottomrule
  \end{tabular}

  \vspace{0.4em}
  {\footnotesize \exitzero{} valid \quad \exitone{} issues found \quad
  \exittwo{} could not run}
\end{frame}

\begin{frame}{What it actually checks}
  \begin{enumerate}
    \item \textbf{Missing} — declared in the template, absent from \file{.env}
    \item \textbf{Placeholders} — \texttt{CHANGE\_ME}, \texttt{YOUR\_KEY\_HERE},
          \texttt{your\_*\_value}, \texttt{TODO}, \texttt{xxx}, \texttt{example}
    \item \textbf{Weak credentials} — short or predictable values in
          \texttt{*\_SECRET}, \texttt{*\_KEY}, \texttt{*\_PASSWORD}, \texttt{*\_TOKEN}
    \item \textbf{Boolean traps} — \texttt{DEBUG="False"} is a non-empty string
    \item \textbf{Duplicates} — the same key assigned twice
    \item \textbf{Formats} (opt-in) — url, port, email
    \item \textbf{\texttt{localhost} under Docker} — only when a
          \file{docker-compose.yml} or \file{Dockerfile} is present
  \end{enumerate}

  \vspace{0.6em}
  \begin{alertblock}{The boolean trap, in one line of Python}
    \texttt{>>> bool(os.getenv("DEBUG"))  \# DEBUG="False"} \\
    \texttt{True}\quad\textemdash\ you just shipped debug mode to production.
  \end{alertblock}
\end{frame}

\begin{frame}[fragile]{A real run}
  \capture{21-validate.txt}
\end{frame}

\begin{frame}[fragile]{\flag{--fix}: what it will and will not invent}
  \capture{33-validate-fix.txt}

  {\footnotesize It generates secrets from the OS CSPRNG, and normalises booleans.
  It will \emph{not} invent a database name — a placeholder insertion is reported
  as a placeholder, not as a repair.}
\end{frame}

\begin{frame}[fragile]{Machine output: \flag{--format json}}
  \capturepart{22-validate-json.txt}{1}{22}

  {\footnotesize Machine output goes to \textbf{stdout}; the human trailer goes to
  \textbf{stderr}. \texttt{evnx validate --format json 2>/dev/null | jq} is safe.}
\end{frame}

\begin{frame}[fragile]{Machine output: \flag{--format github}}
  \capture{23-validate-github.txt}

  {\footnotesize One annotation per finding, on the line the issue is on, with
  the suggestion after a \texttt{\%0A} escape. \texttt{line=1} appears only for
  a variable that is missing and therefore has no line.}
\end{frame}

\begin{frame}[fragile]{Context-aware: \texttt{localhost} under Docker}
  \capturepart{43-validate-localhost-docker.txt}{1}{18}

  {\footnotesize This check fires \emph{only} when Docker configuration is
  present in the project. \texttt{localhost} is correct on a laptop and wrong
  inside a container, so the same value is a finding in one project and not in
  another.}
\end{frame}

\begin{frame}{Use cases}
  \begin{tabular}{@{}p{0.42\textwidth}p{0.52\textwidth}@{}}
    \toprule
    \textbf{Situation} & \textbf{Command} \\
    \midrule
    Pre-flight before \texttt{npm run dev} & \texttt{evnx validate} \\
    Repair the easy problems & \texttt{evnx validate --fix} \\
    Check production config & \texttt{evnx validate --env-name production} \\
    CI gate with annotations & \texttt{evnx validate --format github} \\
    Advisory only, never fail & \texttt{evnx validate --exit-zero} \\
    Suppress a known issue type & \texttt{evnx validate --ignore invalid\_url} \\
    \bottomrule
  \end{tabular}

  \vspace{0.8em}
  \begin{block}{Gate on non-zero, never on \texttt{== 1}}
    \exitone{} means ``I have a verdict and it is bad''. \exittwo{} means
    ``I have \emph{no} verdict''. A pipeline testing \texttt{\$? -eq 1} treats a
    crash as a pass.
  \end{block}
\end{frame}

\section{evnx scan}

\begin{frame}{\cmd{evnx scan} — find secrets that look real}
  \textbf{Purpose.} Detect credentials that appear to be live, in every
  \file{.env} variant in a tree, and say where to revoke them.

  \begin{block}{Synopsis}\texttt{evnx scan [OPTIONS] [PATH]...}\end{block}

  \begin{tabular}{@{}ll@{}}
    \toprule
    \textbf{Flag} & \textbf{What it does} \\
    \midrule
    \texttt{[PATH]...}            & \textbf{Positional}, default \texttt{.}; already recursive \\
    \flag{--severity <LEVEL>}     & \texttt{high} / \texttt{medium} / \texttt{low} (default) \\
    \flag{--pattern <REGEX>}      & Your own key shapes; repeatable \\
    \flag{--exclude <GLOB>}       & Default \file{.env.example} \\
    \flag{--ignore-placeholders}  & Skip values that look like placeholders \\
    \flag{--format <FORMAT>}      & \texttt{pretty}, \texttt{json}, \texttt{sarif}, \texttt{github} \\
    \flag{--exit-zero}            & Suppress \exitone{} only, never \exittwo \\
    \bottomrule
  \end{tabular}

  \vspace{0.4em}
  {\footnotesize \exitzero{} clean \quad \exitone{} secrets found \quad
  \exittwo{} the scan could not be completed}

  \vspace{0.3em}
  \begin{alertblock}{There is no \flag{--path}}
    Paths are positional. \texttt{evnx scan --path .} exits \exittwo.
  \end{alertblock}
\end{frame}

\begin{frame}[fragile]{A real run}
  \capture{24-scan.txt}
\end{frame}

\begin{frame}{Two kinds of finding — and the difference matters}
  \begin{columns}[T]
    \begin{column}{0.48\textwidth}
      \textbf{Matched by \emph{value}}\\[0.3em]
      {\footnotesize The value matches a known provider format. This is the
      line that means \emph{drop everything and rotate}.}
      \vspace{0.4em}
      {\scriptsize\ttfamily
      ! matches a live key format,\\ \phantom{! }not a placeholder\\
      $\rightarrow$ revoke at dashboard.stripe.com}
    \end{column}
    \begin{column}{0.48\textwidth}
      \textbf{Matched by \emph{name}}\\[0.3em]
      {\footnotesize The variable is named like a secret. The value was never
      read, so this is a prompt to look, not to rotate.}
      \vspace{0.4em}
      {\scriptsize\ttfamily
      \textperiodcentered\ flagged by variable name\\ \phantom{. }-- the value was not checked}
    \end{column}
  \end{columns}

  \vspace{1em}
  {\footnotesize In \flag{--format json} this is the \texttt{matched\_by} field
  (\texttt{"name"} or \texttt{"value"}). Attaching ``rotate this now'' to a name
  match is how people learn to ignore a scanner.}
\end{frame}

\begin{frame}[fragile]{\flag{--severity high} — the blocking gate}
  \capture{25-scan-severity-high.txt}

  {\footnotesize \flag{--severity} filters what is reported, counted
  \emph{and} exited on — all three describe the same set. The default
  (\texttt{low}) will block on every high-entropy string, including public keys
  that are meant to be public.}
\end{frame}

\begin{frame}[fragile]{\flag{--format github} — inline PR annotations}
  \capture{27-scan-github.txt}
\end{frame}

\begin{frame}[fragile]{\flag{--format sarif} — the Security tab}
  \capturepart{28-scan-sarif.txt}{1}{20}
\end{frame}

\begin{frame}{What the SARIF carries}
  \begin{itemize}
    \item \textbf{Stable \texttt{ruleId}s} — \texttt{secret/sensitive-config-key},
          not one rule per variable name, so the Security tab does not grow a
          rule per secret.
    \item \textbf{\texttt{rules[]} with \texttt{helpUri}} pointing at
          \texttt{evnx.dev/guides/commands/scan} — the fix is one click from
          the alert.
    \item \textbf{\texttt{partialFingerprints}} keyed on
          \texttt{file::variable::rule}, \emph{not} line number — so inserting a
          line above a secret does not close and reopen a triaged alert.
  \end{itemize}

  \vspace{0.6em}
  \begin{block}{Values are always masked}
    Every format prints first-4/last-4 only. A secret scanner that prints
    secrets into your CI log has not helped you.
  \end{block}
\end{frame}

\begin{frame}{Use cases}
  \begin{tabular}{@{}p{0.40\textwidth}p{0.54\textwidth}@{}}
    \toprule
    \textbf{Situation} & \textbf{Command} \\
    \midrule
    Pre-commit hook & \texttt{evnx scan --severity high} \\
    Block a pull request & \texttt{evnx scan --severity high --format github} \\
    Upload to Security tab & \texttt{evnx scan --format sarif --exit-zero > x.sarif} \\
    Your own token shape & \texttt{evnx scan --pattern 'ACME-[A-Z0-9]\{32\}'} \\
    One directory only & \texttt{evnx scan ./services/api} \\
    \bottomrule
  \end{tabular}

  \vspace{0.8em}
  \begin{alertblock}{A regex that does not compile exits \exittwo}
    And \flag{--exit-zero} does \emph{not} suppress it. A broken rule reading as
    ``no secrets found'' would be CI passing a search that never ran.
  \end{alertblock}
\end{frame}

\section{evnx diff}

\begin{frame}{\cmd{evnx diff} — two commands sharing one name}
  \textbf{Purpose.} Show what differs between two environment files.

  \begin{block}{Synopsis}\texttt{evnx diff [OPTIONS]}\end{block}

  \begin{tabular}{@{}ll@{}}
    \toprule
    \textbf{Flag} & \textbf{What it does} \\
    \midrule
    \flag{--env-name <NAME>}   & Left side: \file{.env.<NAME>} \\
    \flag{--against <NAME>}    & \textbf{Right side}: \file{.env.<NAME>} not the template \\
    \flag{--env} / \flag{--example} & Explicit paths \\
    \flag{--show-values}       & Reveal values (masked by default) \\
    \flag{--ignore-keys <K,K>} & Skip keys that are meant to differ \\
    \flag{--format <FORMAT>}   & \texttt{pretty}, \texttt{json}, patch \\
    \flag{--with-stats}        & Extended statistics in JSON \\
    \flag{--reverse}           & Swap direction \\
    \flag{-i, --interactive}   & Interactive merge (patch format) \\
    \bottomrule
  \end{tabular}
\end{frame}

\begin{frame}[fragile]{Default: against the template}
  \capture{29-diff.txt}

  \begin{alertblock}{This form is not a CI gate}
    \file{.env.example} holds empty placeholders by construction, so every
    filled-in variable counts as a difference. A perfectly healthy project
    exits \exitone.
  \end{alertblock}
\end{frame}

\begin{frame}[fragile]{\flag{--against}: between two real environments}
  \capture{58-diff-against.txt}

  {\footnotesize \textbf{This form \emph{is} a gate} — it exits \exitzero{} when
  the two environments match.}
\end{frame}

\begin{frame}{Which form to use}
  \begin{tabular}{@{}lll@{}}
    \toprule
    \textbf{Invocation} & \textbf{Gateable} & \textbf{Why} \\
    \midrule
    \texttt{evnx diff} & \textcolor{evnxErr}{No} & Template values differ by construction \\
    \texttt{evnx diff --against Y} & \textcolor{evnxOk}{Yes} & Two real files; \exitzero{} when equal \\
    \texttt{evnx sync --check} & \textcolor{evnxOk}{Yes} & Purpose-built, documents 0/1/2 \\
    \bottomrule
  \end{tabular}

  \vspace{1em}
  \textbf{Use cases}
  \begin{itemize}\footnotesize
    \item ``What is staging missing that production has?''\\
          \texttt{evnx diff --env-name production --against staging}
    \item ``Ignore the keys that are supposed to differ''\\
          \texttt{... --ignore-keys APP\_ENV,RUST\_LOG,SENTRY\_ENVIRONMENT}
    \item ``Give me something to diff in CI''\\
          \texttt{... --format json --with-stats}
  \end{itemize}
\end{frame}

\section{evnx sync}

\begin{frame}{\cmd{evnx sync} — keep \file{.env.example} honest}
  \textbf{Purpose.} Propagate variable \emph{names} between \file{.env} and
  \file{.env.example}, never the values.

  \begin{block}{Synopsis}\texttt{evnx sync [OPTIONS]}\end{block}

  \begin{tabular}{@{}ll@{}}
    \toprule
    \textbf{Flag} & \textbf{What it does} \\
    \midrule
    \flag{--direction forward} & \file{.env} $\rightarrow$ \file{.env.example} (default) \\
    \flag{--direction reverse} & \file{.env.example} $\rightarrow$ \file{.env} (onboarding) \\
    \flag{--check}             & \textbf{The CI gate.} Implies \flag{--dry-run} \\
    \flag{-n, --dry-run}       & Preview, write nothing \\
    \flag{-f, --force}         & Skip confirmation — \textbf{required in CI} \\
    \flag{--placeholder}       & Placeholder values for new variables \\
    \flag{--format <FORMAT>}   & \texttt{json} carries a diff summary \\
    \flag{--naming-policy}     & \texttt{warn} (default) / \texttt{error} / \texttt{ignore} \\
    \bottomrule
  \end{tabular}
\end{frame}

\begin{frame}[fragile]{\flag{--check} — the only flag that documents its own exit codes}
\begin{lstlisting}[style=evnxcode]
      --check
          Report whether the files are in step through the exit code.

          Implies --dry-run: nothing is ever written. This is the CI gate.

            0  in sync
            1  out of sync - commit the result of `evnx sync`
            2  error - a file is missing, or will not parse
\end{lstlisting}
  \capture{35-sync-check.txt}
\end{frame}

\begin{frame}[fragile]{Fixing the drift}
  \capture{36-sync-forward.txt}

  \begin{alertblock}{\flag{--force} is required in any non-interactive context}
    Without it \cmd{sync} prompts, and with no TTY exits \exittwo{} with
    \texttt{not a terminal}. Piping \texttt{yes} does not help — the prompt needs
    a real terminal.
  \end{alertblock}
\end{frame}

\begin{frame}{Use cases}
  \begin{tabular}{@{}p{0.42\textwidth}p{0.52\textwidth}@{}}
    \toprule
    \textbf{Situation} & \textbf{Command} \\
    \midrule
    CI: has the template drifted? & \texttt{evnx sync --check} \\
    Commit the fix & \texttt{evnx sync --direction forward --force} \\
    New teammate fills \file{.env} & \texttt{evnx sync --direction reverse --force} \\
    Preview without writing & \texttt{evnx sync --dry-run} \\
    Machine-readable drift & \texttt{evnx sync --check --format json} \\
    \bottomrule
  \end{tabular}

  \vspace{0.8em}
  \begin{block}{Why this matters more than it looks}
    \file{.env.example} is the only contract a new developer has. A variable
    added to \file{.env} and not to the template is a broken checkout for the
    next person — and nothing else in the toolchain notices.
  \end{block}
\end{frame}

\section{evnx convert}

\begin{frame}{\cmd{evnx convert} — render \file{.env} as something else}
  \begin{block}{Synopsis}\texttt{evnx convert [OPTIONS]}\end{block}

  \begin{tabular}{@{}ll@{}}
    \toprule
    \flag{--to <FORMAT>}      & Target; omit for interactive selection \\
    \flag{-o, --output <FILE>} & Or just redirect — stdout is clean \\
    \flag{--include <GLOB>}   & \texttt{"DB\_*,AWS\_*"} \\
    \flag{--exclude <GLOB>}   & Applied \emph{after} \flag{--include} \\
    \flag{--prefix <P>}       & Prepend to every name \\
    \flag{--transform <MODE>} & \texttt{uppercase}, \texttt{lowercase}, \texttt{camelCase}, \texttt{snake\_case} \\
    \flag{--base64}           & Base64-encode values \\
    \flag{--env-name <NAME>}  & Source \file{.env.<NAME>} \\
    \bottomrule
  \end{tabular}

  \vspace{0.5em}
  \textbf{14 formats.}
  {\footnotesize \texttt{json yaml shell aws-secrets gcp-secrets azure-keyvault
  github-actions docker-compose kubernetes terraform doppler heroku vercel
  railway}}

  \vspace{0.3em}
  {\footnotesize \textbf{Aliases:} \texttt{k8s}, \texttt{tf}, \texttt{yml},
  \texttt{gh-actions}, \texttt{compose}}
\end{frame}

\begin{frame}[fragile]{JSON, and why redirection is safe}
  \capture{37-convert-json.txt}

  {\footnotesize The \texttt{Docs:} trailer is on \textbf{stderr}.
  \texttt{evnx convert --to json > config.json} produces clean JSON.}
\end{frame}

\begin{frame}[fragile]{Kubernetes: \texttt{stringData} vs \texttt{data}}
  \begin{columns}[T]
    \begin{column}{0.5\textwidth}
      {\scriptsize without \flag{--base64}}
      \capturepart{38-convert-k8s.txt}{1}{12}
    \end{column}
    \begin{column}{0.5\textwidth}
      {\scriptsize with \flag{--base64}}
      \capturepart{39-convert-k8s-b64.txt}{1}{12}
    \end{column}
  \end{columns}

  {\footnotesize \texttt{stringData} is plaintext in the manifest — Kubernetes
  encodes it itself on apply. Use \flag{--base64} only if you want a
  pre-encoded \texttt{data:} block.}
\end{frame}

\begin{frame}[fragile]{Filtering — least privilege for a build step}
  \capture{41-convert-filtered.txt}

  {\footnotesize A front-end build that needs only \texttt{NEXT\_PUBLIC\_*} has
  no reason to be handed your database password.}
\end{frame}

\begin{frame}{\cmd{convert} vs \cmd{migrate}}
  \begin{tabular}{@{}lll@{}}
    \toprule
     & \cmd{evnx convert} & \cmd{evnx migrate} \\
    \midrule
    Produces & A file or stdout & The platform CLI's own commands \\
    Uploads  & Never & Only \texttt{--to github-actions} \\
    Use for  & k8s, Terraform, compose, IaC & A guided one-time move \\
    \bottomrule
  \end{tabular}

  \vspace{1em}
  \textbf{Use cases}
  \begin{itemize}\footnotesize
    \item \texttt{evnx convert --to kubernetes --include 'DB\_*' -o secret.yaml}
    \item \texttt{evnx convert --to terraform --prefix APP\_ --transform uppercase}
    \item \texttt{evnx convert --env-name production --to docker-compose}
    \item \texttt{evnx convert --to shell > /tmp/env.sh \&\& source /tmp/env.sh}
  \end{itemize}
\end{frame}

\section{evnx template}

\begin{frame}{\cmd{evnx template} — generate a config file from \file{.env}}
  \begin{block}{Synopsis}
    \texttt{evnx template --input <FILE> --output <FILE> [OPTIONS]}
  \end{block}

  \begin{tabular}{@{}ll@{}}
    \toprule
    \flag{--input <FILE>}   & The template (required) \\
    \flag{--output <FILE>}  & Where to write (required) \\
    \flag{--strict}         & Refuse to write with an unresolved placeholder \\
    \flag{--gitignore}      & Add the output to \file{.gitignore}, no prompt \\
    \flag{--no-gitignore}   & Skip gitignore handling entirely \\
    \flag{--env-name <N>}   & Take values from \file{.env.<N>} \\
    \bottomrule
  \end{tabular}

  \vspace{0.6em}
  \textbf{Placeholder syntax}\quad
  \texttt{\{\{NAME\}\}} \quad and \quad \texttt{\{\{NAME|default:value\}\}}
\end{frame}

\begin{frame}[fragile]{The template}
  \capture{50-template-input.txt}
  \vspace{-0.4em}
  \capture{51-template-loose.txt}
\end{frame}

\begin{frame}[fragile]{What it wrote — note the unresolved placeholder}
  \capture{51-template-output.txt}

  \begin{alertblock}{The default is permissive}
    \texttt{\{\{OTEL\_ENDPOINT\}\}} was copied through \emph{literally} and the
    command exited \exitzero. That is how a placeholder ends up in a production
    config as a literal string.
  \end{alertblock}
\end{frame}

\begin{frame}[fragile]{\flag{--strict} — refuse, and say what was not done}
  \capture{52-template-strict.txt}

  {\footnotesize Three things this error does that most do not: it \textbf{names
  the variable}, it \textbf{offers the remedy syntax}, and it states
  \textbf{``Nothing was written''}. Exit \exittwo{} — a refusal to act, not a
  finding.}
\end{frame}

\begin{frame}{Use cases}
  \begin{itemize}
    \item \textbf{nginx / Caddy config from \file{.env}}\\
      {\footnotesize\texttt{evnx template --input nginx.conf.tpl --output nginx.conf --strict}}
    \item \textbf{Per-environment app config}\\
      {\footnotesize\texttt{evnx template --input app.yaml.tpl --output app.yaml --env-name production --strict}}
    \item \textbf{Optional settings that must not break \flag{--strict}}\\
      {\footnotesize\texttt{pool: \{\{DB\_POOL\_SIZE|default:10\}\}}}
  \end{itemize}

  \vspace{1em}
  \begin{block}{Always \flag{--strict} in CI}
    Without it, a missing variable becomes a literal string in a generated
    config and the pipeline goes green.
  \end{block}
\end{frame}

\section{evnx migrate}

\begin{frame}{\cmd{evnx migrate} — move secrets to a manager}
  \begin{block}{Synopsis}\texttt{evnx migrate --to <DEST> [OPTIONS]}\end{block}

  \textbf{Nine destinations.}
  {\footnotesize \texttt{github-actions}, \texttt{aws-secrets-manager},
  \texttt{doppler}, \texttt{infisical}, \texttt{gcp-secret-manager},
  \texttt{azure-keyvault}, \texttt{vercel}, \texttt{heroku}, \texttt{railway}}

  \vspace{0.5em}
  \begin{alertblock}{Only \texttt{--to github-actions} uploads}
    The other eight \textbf{print the platform's own CLI commands} for you to
    run. This is by design and the summary now says so.
  \end{alertblock}

  \vspace{0.3em}
  \begin{tabular}{@{}ll@{}}
    \flag{--dry-run}          & Preview; runs headless \\
    \flag{--include/--exclude} & Comma-separated globs \\
    \flag{--strip-prefix/--add-prefix} & Rewrite names on the way \\
    \flag{--skip-existing/--overwrite} & \texttt{github-actions} only \\
  \end{tabular}
\end{frame}

\begin{frame}[fragile]{A real dry run}
  \capture{43-migrate-doppler.txt}
\end{frame}

\begin{frame}[fragile]{AWS Secrets Manager}
  \capturepart{42-migrate-aws.txt}{1}{22}
\end{frame}

\begin{frame}[fragile]{It refuses rather than guessing}
  \capture{44-migrate-noto.txt}

  {\footnotesize Before v0.5.0, \cmd{migrate} with no \flag{--to} and no
  terminal silently selected the \emph{first} destination in the list, printed a
  plan for somewhere you never chose, and exited \exitzero.}
\end{frame}

\begin{frame}{Use cases}
  \begin{itemize}
    \item \textbf{Seed a GitHub repo's Actions secrets} (the one that uploads)\\
      {\footnotesize\texttt{evnx migrate --to github-actions --repo owner/repo}}
    \item \textbf{Move production config into AWS}\\
      {\footnotesize\texttt{evnx migrate --to aws-secrets-manager --secret-name prod/api/config --dry-run}}
    \item \textbf{Only the database variables}\\
      {\footnotesize\texttt{evnx migrate --to doppler --project app --include 'DB\_*'}}
  \end{itemize}

  \vspace{0.8em}
  \begin{block}{Migration is once per project}
    For day-to-day secret delivery, use \cmd{evnx cloud run} instead — see the
    use-case deck.
  \end{block}
\end{frame}

\section{evnx doctor}

\begin{frame}{\cmd{evnx doctor} — is this project set up correctly?}
  \textbf{Purpose.} Check the things around \file{.env}: gitignore coverage,
  file permissions, syntax, project structure.

  \begin{block}{Synopsis}\texttt{evnx doctor [OPTIONS]}\end{block}

  \begin{tabular}{@{}ll@{}}
    \toprule
    \flag{--fix}             & Repair gitignore coverage and permissions \\
    \flag{--strict}          & Treat warnings as failures \\
    \flag{--path <PATH>}     & Diagnose another directory \\
    \flag{--format <FORMAT>} & \texttt{json} for CI \\
    \bottomrule
  \end{tabular}

  \vspace{0.5em}
  {\footnotesize \texttt{EVNX\_OUTPUT\_JSON=1} is also honoured, and is what the
  older guides document.}
\end{frame}

\begin{frame}[fragile]{A real run}
  \capture{30-doctor.txt}
\end{frame}

\begin{frame}{What it checks, and why each matters}
  \begin{tabular}{@{}p{0.36\textwidth}p{0.58\textwidth}@{}}
    \toprule
    \textbf{Check} & \textbf{Why} \\
    \midrule
    \file{.env} in \file{.gitignore} & The single highest-value check here \\
    File permissions & \texttt{0600}; \texttt{0644} exposes secrets to other users \\
    \file{.env} syntax & Reports \emph{every} bad line, not just the first \\
    \file{.env.example} exists \& tracked & The contract new developers rely on \\
    Project type detection & Confirms evnx understands your stack \\
    Docker configuration & Detects \file{docker-compose.yml} / \file{Dockerfile} \\
    \bottomrule
  \end{tabular}

  \vspace{0.8em}
  {\footnotesize \textbf{v0.5.0:} \cmd{doctor} and the core parser now agree.
  \texttt{export FOO=bar} is valid (it is standard dotenv), a leading underscore
  is valid, and all syntax errors are reported in one run.}
\end{frame}

\begin{frame}{Use cases}
  \begin{itemize}
    \item \textbf{First thing in a cloned repo} — \texttt{evnx doctor}
    \item \textbf{Fix the mechanical problems} — \texttt{evnx doctor --fix}
    \item \textbf{A monorepo service} — \texttt{evnx doctor --path services/api}
    \item \textbf{CI health check} — \texttt{evnx doctor --format json | jq '.summary.errors'}
    \item \textbf{Strict gate} — \texttt{evnx doctor --strict}
  \end{itemize}

  \vspace{1em}
  \begin{block}{\cmd{doctor} vs \cmd{validate}}
    \cmd{doctor} asks \emph{``is this project set up correctly?''} —
    gitignore, permissions, structure.\\
    \cmd{validate} asks \emph{``are the values in this file fit to run?''} —
    missing, placeholder, weak, malformed.
  \end{block}
\end{frame}

\section{evnx spec}

\begin{frame}{\cmd{evnx spec} — the contract \file{.env.example} cannot express}
  \textbf{Purpose.} Declare what each variable \emph{is}: required, secret,
  a URL, a port, production-only.

  \begin{block}{Synopsis}\texttt{evnx spec init}\end{block}

  \file{.env.example} can say a variable exists. It cannot say it is required
  only in production, must parse as a URL, or holds a credential despite a name
  that gives nothing away.

  \vspace{0.6em}
  \begin{block}{Where it lives}
    \texttt{[vars.*]} tables in \file{.evnx.toml} — a committed, project-level
    file. Read by \cmd{validate}, \cmd{scan} and \cmd{sync}.
  \end{block}
\end{frame}

\begin{frame}[fragile]{\cmd{spec init} — infer a first draft}
  \capture{31-spec-init.txt}
\end{frame}

\begin{frame}[fragile]{The generated \file{.evnx.toml}}
  \capturepart{31-spec-toml.txt}{1}{24}

  {\footnotesize The header is honest: \emph{``Every field below is inferred
  from this project's files — review it, then commit it.''} Output is
  alphabetically sorted and deterministic.}
\end{frame}

\begin{frame}[fragile]{Hand-written entries are where the value is}
\begin{lstlisting}[style=evnxcode]
[vars.TENANT_A_TOKEN]
secret = true                    # scan reports it; no heuristic could
environments = ["production"]    # required only in production

[vars.PORT]
format = "port"
description = "HTTP listen port"
\end{lstlisting}

  \begin{block}{The asymmetry that matters}
    \texttt{secret = true} \textbf{adds} a finding \cmd{scan} could not otherwise
    reach. \texttt{secret = false} \textbf{retracts} a guess made from a
    variable's \emph{name} — but never a match on its \emph{value}. A variable
    holding \texttt{sk\_live\_\ldots} is reported however it is declared, because
    \file{.evnx.toml} is committed and one wrong line would otherwise silence a
    live key for everyone who clones the repository.
  \end{block}
\end{frame}

\begin{frame}[fragile]{Settings that weaken scanning are announced}
\begin{lstlisting}[style=evnxcode]
  config    .evnx.toml (scan.severity=high, vars.secret=false on 1 variable(s))
\end{lstlisting}

  {\footnotesize Because the file is committed, one commit could otherwise
  narrow scanning for everyone who clones the repository with nothing on the
  command line to show it.}

  \vspace{0.8em}
  \textbf{Use cases}
  \begin{itemize}\footnotesize
    \item A token whose name reveals nothing: \texttt{secret = true}
    \item A variable only production needs: \texttt{environments = ["production"]}
    \item A shared team scan rule: \texttt{[[scan.patterns]]}
    \item Project-wide severity floor: \texttt{[scan] severity = "high"}
  \end{itemize}
\end{frame}

\section{evnx backup / restore}

\begin{frame}{\cmd{evnx backup} \& \cmd{evnx restore}}
  \textbf{Purpose.} An encrypted copy of a \file{.env} you can store anywhere.

  \vspace{0.4em}
  \begin{columns}[T]
    \begin{column}{0.49\textwidth}
      \texttt{evnx backup [ENV] [OPTS]}
      \begin{tabular}{@{}l@{}}
        \flag{--output <FILE>} \\
        \flag{--key-file <PATH>} \\
        \flag{--keep <N>} \\
        \flag{--env-name <NAME>} \\
        \flag{--verify} \\
      \end{tabular}
    \end{column}
    \begin{column}{0.49\textwidth}
      \texttt{evnx restore <BACKUP> [OPTS]}
      \begin{tabular}{@{}l@{}}
        \flag{--output <FILE>} \\
        \flag{--password-file <PATH>} \\
        \flag{--inspect} \\
        \flag{--dry-run} \\
      \end{tabular}
    \end{column}
  \end{columns}

  \vspace{0.6em}
  \begin{block}{Each side also accepts the other's spelling}
    \cmd{backup} documents \flag{--key-file} and \cmd{restore} documents
    \flag{--password-file}, but \emph{both} accept \emph{both} — verified — and
    both honour \texttt{EVNX\_PASSWORD}. One key file always round-trips,
    including a binary one.
  \end{block}

  \vspace{0.3em}
  {\footnotesize Crypto: Argon2id over the passphrase $\rightarrow$ AES-256-GCM.
  Output is Base64.}
\end{frame}

\begin{frame}[fragile]{Backing up}
  \capture{53-backup.txt}
\end{frame}

\begin{frame}[fragile]{\flag{--verify}: decrypt it before trusting it}
  \capture{53b-backup-verify.txt}

  {\footnotesize Without \flag{--verify} the summary reports
  \texttt{Verified: no} — the backup was written but never read back. The advice
  ``test the restore before deleting the original'' is only as good as the test
  you actually run.}
\end{frame}

\begin{frame}[fragile]{\flag{--inspect}: what is in this backup?}
  \capture{54-restore-inspect.txt}

  {\footnotesize Key names only. Values are never printed, nothing is written,
  and no overwrite prompt appears. This is the right way to check a backup
  before committing to it.}
\end{frame}

\begin{frame}[fragile]{The round trip, verified}
  \capture{56-restore.txt}
  \vspace{-0.3em}
  \capture{57-restore-diff.txt}
\end{frame}

\begin{frame}{Exit codes are richer here — deliberately}
  \begin{tabular}{@{}cl@{}}
    \toprule
    \textbf{Code} & \textbf{Meaning} \\
    \midrule
    \exitzero & Success \\
    \exittwo  & Wrong password or corrupt backup \\
    \texttt{3} & Backup file not found \\
    \texttt{4} & Cancelled (informational) \\
    \texttt{5} & Validation fallback (informational) \\
    \texttt{6} & Integrity check failed \\
    \bottomrule
  \end{tabular}

  \vspace{0.8em}
  \textbf{Use cases}
  \begin{itemize}\footnotesize
    \item Nightly, with rotation: \texttt{evnx backup --env-name production --key-file /run/secrets/k --keep 7}
    \item CI restore with no file on disk: \texttt{EVNX\_PASSWORD="\$P" evnx restore prod.backup}
    \item Before a risky edit: \texttt{evnx backup \&\& \ldots}
  \end{itemize}

  \vspace{0.4em}
  \begin{alertblock}{There is deliberately no \flag{--force} on restore}
    Restoring over an existing file always prompts. The passphrase cannot be
    recovered — if you lose it, the backup is gone.
  \end{alertblock}
\end{frame}

\section{evnx auth}

\begin{frame}{\cmd{evnx auth} — your evnx cloud account}
  \begin{block}{Synopsis}\texttt{evnx auth <SUBCOMMAND>}\end{block}

  \begin{tabular}{@{}ll@{}}
    \toprule
    \texttt{auth register} & Create an account (keys generated locally) \\
    \texttt{auth login}    & SRP-6a login, TOTP challenge if enabled \\
    \texttt{auth logout}   & End the session \\
    \texttt{auth status}   & Who am I, against which server \\
    \midrule
    \texttt{auth token create/list/revoke} & CI tokens (\texttt{evnx\_tok\_\ldots}) \\
    \texttt{auth totp enable/disable}      & Second factor \\
    \texttt{auth totp recovery-codes}      & 10 single-use codes, shown once \\
    \texttt{auth sessions list/revoke}     & Active sessions \\
    \texttt{auth sessions revoke-others}   & After a laptop goes missing \\
    \bottomrule
  \end{tabular}

  \vspace{0.5em}
  {\footnotesize Credentials live in \file{\textasciitilde/.config/evnx/},
  mode \texttt{0600} — the \file{\textasciitilde/.aws/credentials} model.}
\end{frame}

\begin{frame}[fragile]{The login flow — your password never leaves the device}
\begin{semiverbatim}\scriptsize
POST /auth/register    SRP verifier + salts + public keys (all client-side)
GET  /auth/verify-email   emailed link flips email_verified
POST /auth/srp/init    -> session_id, salts, server public B
POST /auth/srp/verify  M1 -> M2 + \{ access, refresh \}
                          -> M2 + totp_pending_token   (if TOTP on)
POST /auth/totp/verify pending token + code -> \{ access, refresh \}
\end{semiverbatim}

  \vspace{0.4em}
  \begin{block}{SRP-6a in one sentence}
    The server stores a \emph{verifier}, not a password hash it could brute
    force offline in the usual way, and the proof exchange never transmits the
    password or anything derived from it that a passive observer could reuse.
  \end{block}

  {\footnotesize 15-minute access token; 30-day refresh token, rotated on use.}
\end{frame}

\begin{frame}[fragile]{The transport guard}
  \capture{66-tls-guard.txt}

  {\footnotesize Loopback is exempt, so local development against
  \texttt{http://127.0.0.1:8080} works. The error explains the \emph{threat},
  not just the rule.}
\end{frame}

\begin{frame}{Use cases}
  \begin{itemize}\footnotesize
    \item \textbf{CI token, scoped to one vault, read-only}\\
      \texttt{evnx auth token create --scope read\_only --vault app/production}
    \item \textbf{Turn on 2FA} — \texttt{evnx auth totp enable} (QR in the terminal)
    \item \textbf{Save the recovery codes} — \texttt{evnx auth totp recovery-codes}
    \item \textbf{Lost laptop} — \texttt{evnx auth sessions revoke-others}
  \end{itemize}

  \vspace{0.8em}
  \begin{block}{Minting a token requires a real login}
    \texttt{/auth/tokens} sits behind a user-session guard, so a leaked CI token
    cannot mint itself a longer-lived, wider-scoped replacement and survive
    revocation of the original.
  \end{block}
\end{frame}

\section{evnx vault}

\begin{frame}{\cmd{evnx vault} — the unit of sharing}
  \begin{block}{Synopsis}\texttt{evnx vault <SUBCOMMAND>}\end{block}

  \begin{tabular}{@{}ll@{}}
    \toprule
    \texttt{vault create <NAME>} & \texttt{app} or \texttt{app/production} \\
    \texttt{vault list}          & Vaults you can reach \\
    \texttt{vault delete <NAME>} & Soft delete \\
    \texttt{vault share <NAME> --with <EMAIL>} & Wrap the key for someone else \\
    \texttt{vault members <NAME>} & Who has access, at what role \\
    \texttt{vault role <NAME> --user --role} & Change a role \\
    \texttt{vault revoke <NAME> --user <EMAIL>} & Remove \textbf{and re-key} \\
    \bottomrule
  \end{tabular}

  \vspace{0.4em}
  \begin{alertblock}{Two names for the same thing}
    \cmd{vault share} takes \flag{--with}; \cmd{vault revoke} and
    \cmd{vault role} take \flag{--user}. Same concept, two flags --- worth
    knowing before you guess.
  \end{alertblock}

  \vspace{0.3em}
  {\footnotesize A vault name is \texttt{name} or \texttt{name/environment};
  \texttt{app/production} binds naturally to \file{.env.production}.}
\end{frame}

\begin{frame}{Two wrap modes — and the distinction is load-bearing}
  \begin{tabular}{@{}p{0.26\textwidth}p{0.40\textwidth}p{0.24\textwidth}@{}}
    \toprule
    \textbf{Whose copy} & \textbf{How it is wrapped} & \textbf{Post-quantum} \\
    \midrule
    The creator's own & XChaCha20 under an HKDF subkey of the master key & \evnxCheck{} symmetric only \\
    A member's, shared & \textbf{hybrid X25519 + ML-KEM-768} & \evnxCheck{} holds if either survives \\
    \bottomrule
  \end{tabular}

  \vspace{0.8em}
  The creator's copy deliberately does \emph{not} go through ECDH — that would
  put every solo vault behind X25519 and expose it to
  harvest-now-decrypt-later.

  \vspace{0.4em}
  \begin{block}{There is no classical-only path}
    Not in the crate, not in the CLI, and not in the database, where a CHECK
    constraint makes half a wrap unstorable. Sharing with an account that has no
    ML-KEM key on file is \textbf{refused}, not downgraded.
  \end{block}
\end{frame}

\begin{frame}[fragile]{Revocation actually revokes}
  \capturepart{63-vault-revoke-help.txt}{1}{8}

  {\footnotesize This is the best-written help in the CLI because it states the
  \emph{limit} of what the command can do — the sentence most security tooling
  omits.}
\end{frame}

\begin{frame}{The honest caveat}
  \begin{alertblock}{The recipient's public keys come from the server}
    A malicious server could substitute its own and read what you share. There
    is no third party to check against, and out-of-band fingerprint verification
    is not built.
  \end{alertblock}

  \vspace{0.6em}
  So \emph{``the server cannot read your secrets''} becomes \emph{``\ldots cannot
  read them \textbf{passively}''} the moment you share.

  \vspace{0.8em}
  \textbf{Use cases}
  \begin{itemize}\footnotesize
    \item \texttt{evnx vault create app/production}
    \item \texttt{evnx vault share app/production --with teammate@example.com}
    \item \texttt{evnx vault role app/production --user t@x.com --role viewer}
    \item \texttt{evnx vault revoke app/production --user former@example.com}
  \end{itemize}
\end{frame}

\section{evnx cloud}

\begin{frame}{\cmd{evnx cloud} — encrypted sync}
  \begin{block}{Synopsis}\texttt{evnx cloud <SUBCOMMAND>}\end{block}

  \begin{tabular}{@{}ll@{}}
    \toprule
    \texttt{cloud push}   & Encrypt locally, upload ciphertext \\
    \texttt{cloud pull}   & Download and decrypt \\
    \texttt{cloud run -- <CMD>} & Inject into a subprocess; nothing on disk \\
    \texttt{cloud status} & Server, account, token state \\
    \texttt{cloud history}& Version list \\
    \texttt{cloud link <VAULT>} & Bind this directory to a vault \\
    \texttt{cloud unlink} & Remove the binding \\
    \bottomrule
  \end{tabular}
\end{frame}

\begin{frame}[fragile]{The zero-knowledge chain}
\begin{semiverbatim}\scriptsize
Password
  -> Argon2id(password, salt, m=64MB, t=3, p=4)
       -> MasterKey            (never leaves your device)
           -> XChaCha20(MK) wraps VaultKey   (random 256-bit, one per vault)
               -> AES-256-GCM(VK, nonce) encrypts .env
                   -> ciphertext -> object storage
\end{semiverbatim}

  \vspace{0.5em}
  \begin{block}{Associated data stops a rollback}
    Blobs are encrypted with \texttt{vault\_aad(vault\_id, version)} bound in, so
    a malicious server cannot replay an old version as current.

    \vspace{0.3em}
    {\footnotesize Consequence: a \texttt{409} conflict requires
    \textbf{re-encryption}, not a bare retry — the version number is
    authenticated into the ciphertext.}
  \end{block}
\end{frame}

\begin{frame}[fragile]{\cmd{cloud run} — secrets without a file}
  \capturepart{61-cloud-run-help.txt}{1}{12}
\end{frame}

\begin{frame}{Why \cmd{cloud run} is different}
  \begin{itemize}
    \item Values travel in the \textbf{environment block}, never the argument
          vector — so nothing in \texttt{ps}, nothing in shell history,
          nothing on disk.
    \item \flag{--include} / \flag{--exclude} narrow what the child sees.
          A build step needing only \texttt{NEXT\_PUBLIC\_*} has no reason to
          hold your database password.
    \item A filter matching nothing is an \textbf{error}, not a silent run with
          zero secrets.
    \item The child's exit code comes back unchanged.
    \item It runs the command directly, not through a shell. For a pipeline,
          ask for one: \texttt{-- sh -c 'a | b'}.
  \end{itemize}

  \vspace{0.6em}
  \begin{alertblock}{Be precise about what it removes}
    It removes the plaintext \emph{file}, not the master password. The vault key
    is wrapped under your master key, so decryption needs it wherever this runs.
    In CI, feed it with \flag{--password-stdin}.
  \end{alertblock}
\end{frame}

\begin{frame}[fragile]{\cmd{cloud status}}
  \capture{65-cloud-status.txt}

  {\footnotesize A directory bound with \cmd{cloud link} needs no
  \flag{--vault} on any subsequent command.}
\end{frame}

\section{evnx completions}

\begin{frame}{\cmd{evnx completions} — shell tab-completion}
  \begin{block}{Synopsis}\texttt{evnx completions <SHELL>}\end{block}

  Supported: \texttt{bash} \quad \texttt{zsh} \quad \texttt{fish} \quad
  \texttt{powershell}

  \vspace{0.6em}
  \begin{tabular}{@{}llr@{}}
    \toprule
    \textbf{Shell} & \textbf{Install to} & \textbf{Lines} \\
    \midrule
    bash       & \file{\textasciitilde/.local/share/bash-completion/completions/evnx} & 3{,}432 \\
    zsh        & \file{\$\{fpath[1]\}/\_evnx} (must be named \texttt{\_evnx}) & 2{,}771 \\
    fish       & \file{\textasciitilde/.config/fish/completions/evnx.fish} & 551 \\
    powershell & append to \texttt{\$PROFILE} & 1{,}376 \\
    \bottomrule
  \end{tabular}

  \vspace{0.6em}
  {\footnotesize An unsupported shell is exit \exittwo{} and names the four
  valid values.}
\end{frame}

\begin{frame}[fragile]{Installing}
\begin{lstlisting}[style=evnxcode,language=bash]
# bash - current user
evnx completions bash > ~/.local/share/bash-completion/completions/evnx

# zsh - the file MUST be called _evnx and sit on your $fpath
evnx completions zsh > "${fpath[1]}/_evnx"

# fish
evnx completions fish > ~/.config/fish/completions/evnx.fish

# powershell - append to the profile, not per-session
evnx completions powershell | Out-String | Invoke-Expression
\end{lstlisting}

  \begin{alertblock}{Regenerate after upgrading}
    Completion data is baked in at build time. A v0.4.x completion file does not
    know about \cmd{spec}, \cmd{cloud run}, \flag{--from-source} or
    \flag{--against}.
  \end{alertblock}
\end{frame}


% ---------------------------------------------------------------------------
\section*{Closing}
\begin{frame}{Where to go next}
  \begin{itemize}
    \item \textbf{Guides} — \url{https://www.evnx.dev}
    \item \textbf{Source} — \url{https://github.com/urwithajit9/evnx}
    \item \textbf{Install} — \texttt{curl -fsSL https://dotenv.space/install.sh | sh}
  \end{itemize}

  \vspace{1em}
  \begin{block}{Companion decks}
    \textbf{Use cases} — workflows that combine commands, and the cloud story.\\
    \textbf{Architecture} — how the CLI is built, for contributors.
  \end{block}
\end{frame}

\end{document}
